\documentclass[11pt]{article}

\usepackage[T1]{fontenc}
\usepackage[margin=1in]{geometry}
\usepackage{amsmath}
\usepackage{amssymb}
\usepackage{booktabs}
\usepackage{hyperref}
\usepackage{xcolor}

\hypersetup{
    colorlinks=true,
    linkcolor=blue,
    urlcolor=blue
}

\setlength{\parindent}{0pt}
\setlength{\parskip}{0.65em}

\newcommand{\BF}{\operatorname{BF}}
\newcommand{\Parents}{\operatorname{Parents}}
\newcommand{\Children}{\operatorname{Children}}
\newcommand{\MinMode}{\mathsf{Min}}
\newcommand{\MaxMode}{\mathsf{Max}}
\newcommand{\Strongest}{\operatorname{Strongest}}

\title{Graph Pruning and Bayes-Factor Calculation\\
       Step-by-Step Algorithm and Implementation Notes}
\author{LogicLikely Structural Insights Engine}
\date{\today}

\begin{document}

\maketitle
\tableofcontents
\newpage

\section{Purpose and Scope}

This document describes the two-stage algorithm implemented by
\texttt{GraphBayesFactorPruner} and \texttt{GraphBayesFactorCalculator}.
The stages are deliberately separate:

\begin{enumerate}
    \item The \emph{graph-pruning stage} selects a compatible collection of
          strong evidence-to-hypothesis paths.
    \item The \emph{Bayes-factor stage} evaluates the pruned graph using two
          conditional probabilities on every retained edge.
\end{enumerate}

Neither stage updates posterior odds.  Leaf Bayes factors are supplied
explicitly to the calculator, and the result is an ordinary Bayes factor for
the requested hypothesis.

The complete caller-side flow is therefore
\[
    G_H^{\star}
        = \operatorname{Prune}(G,H),
    \qquad
    \BF_H
        = \operatorname{Calculate}(G_H^{\star},H,\mathcal{B}_{\mathrm{leaf}}),
\]
where \(G_H^{\star}\) is the pruned graph and
\(\mathcal{B}_{\mathrm{leaf}}\) maps relevant leaf-node identifiers to their
externally supplied Bayes factors.

\section{Graph Model and Direction}

Let
\[
    G=(V,E)
\]
be a directed graph.  An edge is directed from a child proposition toward its
parent proposition:
\[
    e=(u,v) \in E
    \quad\Longleftrightarrow\quad
    u \longrightarrow v.
\]

Thus evidence is encountered on the \emph{from} side of an edge, and the
hypothesis is reached by repeatedly following edges in their stored direction.
For an edge \(e=(u,v)\), the implementation stores three distinct quantities:

\begin{center}
\begin{tabular}{lll}
\toprule
Symbol & Implementation field & Meaning \\
\midrule
\(r_e\) & \texttt{ImportanceToParent}
    & positive path-ranking likelihood ratio \\
\(a_e\) & \texttt{ProbabilityGivenParent}
    & \(P(u\mid v)\) \\
\(b_e\) & \texttt{ProbabilityGivenNotParent}
    & \(P(u\mid\neg v)\) \\
\bottomrule
\end{tabular}
\end{center}

The pruning stage currently uses only \(r_e\).  The Bayes-factor stage uses
only \(a_e\) and \(b_e\).  This distinction is important: the nonlinear edge
transform used during Bayes-factor calculation does not, without additional
assumptions, reduce to one static likelihood ratio suitable for path ranking.
The two probabilities are separately specified model parameters; neither is
determined by the other.  The stored edge kind (for example, \texttt{support} or
\texttt{rebut}) does not alter either numerical algorithm.

For a node \(v\), define
\begin{align*}
    \Parents(v)
        &= \{e=(v,u)\in E\},
        &&\text{edges continuing from \(v\) toward a parent},\\
    \Children(v)
        &= \{e=(u,v)\in E\},
        &&\text{edges entering \(v\) from direct children}.
\end{align*}

The selected hypothesis node is denoted by \(H\).  Its target-relevant cone is
\[
    R_H
      = \{v\in V : \text{there is a directed path from \(v\) to \(H\)}\}.
\]
The zero-edge path ensures that \(H\in R_H\).

For that hypothesis, the relevant evidence-source set is
\[
    \mathcal{S}_H
      = \{v\in R_H\setminus\{H\} :
           \operatorname{kind}(v)\text{ is \texttt{evidence} or
           \texttt{objection}}\}.
\]

\section{Assumptions and Preconditions}

The algorithms use the following assumptions.

\begin{enumerate}
    \item Node identifiers and edge identifiers are unique.
    \item Every edge endpoint refers to an existing node.
    \item The portion of the graph that can reach \(H\) is acyclic.
    \item Every relevant path-ranking weight satisfies \(r_e>0\).
    \item Every retained conditional probability satisfies
          \(0\leq a_e\leq 1\) and \(0\leq b_e\leq 1\).
    \item Every relevant leaf has a supplied Bayes factor greater than zero.
    \item Claim co-dependencies have been removed before multiplying child
          contributions.  The product recurrence treats direct child
          contributions as conditionally independent at their parent.
\end{enumerate}

The implementation recognizes node kinds \texttt{evidence} and
\texttt{objection}, case-insensitively, as evidence sources for pruning.
By contrast, a leaf for BF evaluation is defined structurally by the absence
of relevant child edges, regardless of its node kind.

\section{Stage I: Graph Pruning}

\subsection{Desired Output}

The pruning stage returns a new graph containing the selected hypothesis,
initially strongest evidence-path prefixes, and locally strongest compatible
shared suffixes.  It enforces the following structural rule:

\begin{quote}
All prefixes entering a merge are preserved, but a merge has only one
continuation toward the hypothesis.
\end{quote}

Consequently, several evidence branches may enter the same node, but every
retained non-hypothesis node has one retained outgoing continuation toward
\(H\).

\subsection{Step 1: Build Adjacency Indexes and Validate Identifiers}

The pruner first checks for duplicate node and edge identifiers.  The BF
calculator assumes its input was produced by this stage and therefore relies
on the same uniqueness precondition.  The graph context then builds the
following indexes:

\begin{itemize}
    \item a node-identifier-to-node dictionary;
    \item outgoing parent edges keyed by their child, corresponding to
          \(\Parents(v)\);
    \item incoming child edges keyed by their parent, corresponding to
          \(\Children(v)\).
\end{itemize}

These indexes allow both evidence-to-hypothesis and hypothesis-to-evidence
traversals without rescanning all edges.

\subsection{Step 2: Restrict Work to Nodes That Can Reach the Hypothesis}

The target-relevant node set is
\[
    R_H
      = \{v\in V : \text{there is a directed path from \(v\) to \(H\)}\}.
\]

The algorithm starts a queue at \(H\) and follows child edges in reverse graph
direction.  Whenever it sees an edge \((u,v)\) while visiting \(v\), it adds
\(u\) to the queue if \(u\) has not already been discovered.

This step has three practical effects:

\begin{enumerate}
    \item disconnected evidence is ignored;
    \item edges continuing above or away from the selected hypothesis are
          ignored;
    \item cycle and edge-weight validation for pruning are target-local.
\end{enumerate}

Endpoint validation is the exception: the adjacency context rejects a
dangling edge anywhere in the input graph before \(R_H\) is collected, even
if that malformed edge is disconnected from \(H\).

The hypothesis itself is always in \(R_H\).

\subsection{Step 3: Define Log-Likelihood Path Scores}

Each path-ranking edge weight is converted to log space:
\[
    \ell_e = \log r_e.
\]

For a path
\[
    \pi=(e_1,e_2,\ldots,e_t),
\]
its log-likelihood-ratio score is
\[
    L(\pi)
      = \sum_{i=1}^{t}\ell_{e_i}
      = \log\left(\prod_{i=1}^{t}r_{e_i}\right).
\]

Using sums of logarithms avoids repeatedly multiplying path weights and makes
supporting and counter paths symmetric around the neutral value
\(L=0\), which corresponds to a path likelihood ratio of \(1\).

\subsection{Step 4: Calculate the Min/Max Dynamic-Programming Table}

For every \(v\in R_H\), the algorithm stores
\[
    T(v)=
    \bigl(m(v),M(v),p_{\min}(v),p_{\max}(v),d(v)\bigr),
\]
where

\begin{itemize}
    \item \(m(v)\) is the minimum log score of any path from \(v\) to \(H\);
    \item \(M(v)\) is the maximum log score of any path from \(v\) to \(H\);
    \item \(p_{\min}(v)\) is the first edge on the selected minimum path;
    \item \(p_{\max}(v)\) is the first edge on the selected maximum path;
    \item \(d(v)\) is the longest-hop distance from \(v\) to \(H\).
\end{itemize}

The base case is
\[
    m(H)=M(H)=0,
    \qquad
    p_{\min}(H)=p_{\max}(H)=\varnothing,
    \qquad
    d(H)=0.
\]

For \(v\neq H\), consider each relevant continuation edge
\(e=(v,u)\).  Once \(T(u)\) is known, this edge produces candidates
\begin{align*}
    c_{\min}(e) &= \ell_e + m(u),\\
    c_{\max}(e) &= \ell_e + M(u).
\end{align*}

The recurrence is
\begin{align*}
    m(v)
      &= \min_{e=(v,u)}\bigl(\ell_e+m(u)\bigr),\\
    M(v)
      &= \max_{e=(v,u)}\bigl(\ell_e+M(u)\bigr),\\
    d(v)
      &= 1+\max_{e=(v,u)}d(u).
\end{align*}

The corresponding minimizing and maximizing edges are stored in
\(p_{\min}(v)\) and \(p_{\max}(v)\).  Edge identifiers, rather than only
destination node identifiers, are stored because parallel edges are allowed.

If two candidates have exactly the same computed log score, the edge with the
ordinally earlier identifier is selected; the destination identifier is the
secondary deterministic tie-break.  This makes the selected identifier sets
independent of input edge ordering.  The returned node and edge lists still
preserve the relative order of the original graph objects that survive.

\subsection{Step 5: Evaluate the Table in Reverse Topological Order}

For every relevant node, let
\[
    q_{\mathrm{parent}}(v)
      = \left|\{(v,u)\in E : u\in R_H\}\right|
\]
be the number of unresolved continuations from \(v\) toward \(H\).

The algorithm initially queues nodes for which
\(q_{\mathrm{parent}}(v)=0\).  In a valid target-relevant DAG this begins with
\(H\).  After a node's table entry is computed, the algorithm visits its child
edges and decrements the corresponding child's unresolved count.  A child is
queued when all of its possible continuations have already been evaluated.

If fewer than \(|R_H|\) nodes are processed, the target-relevant subgraph
contains a cycle and pruning stops with an error.

\subsection{Step 6: Choose the Strongest Extreme at a Node}

Path strength is defined as distance from the neutral log score zero.  The
locally strongest mode at \(v\) is therefore
\[
    \Strongest(v)
      =
      \begin{cases}
        \MinMode, & |m(v)|>|M(v)|,\\
        \MaxMode, & |m(v)|\leq |M(v)|.
      \end{cases}
\]

Thus a sufficiently strong counter path can defeat a supporting path with a
larger raw log score.  An exact absolute-strength tie selects the maximum mode,
matching the existing maximum-mode tie convention.

Once a mode has been selected, it is followed consistently:

\begin{itemize}
    \item \(\MinMode\) follows \(p_{\min}(v)\) and remains in
          \(\MinMode\) at the next node;
    \item \(\MaxMode\) follows \(p_{\max}(v)\) and remains in
          \(\MaxMode\) at the next node.
\end{itemize}

The mode is changed only when multiple path components merge at a node.

\subsection{Step 7: Create Distance Buckets and Arrival Lists}

If \(\mathcal{S}_H=\varnothing\), the algorithm skips the arrival sweep and
returns a graph containing only \(H\).  Otherwise, let
\[
    D=\max_{s\in\mathcal{S}_H} d(s)
\]
be the largest distance of any relevant evidence node, where
\(\mathcal{S}_H\) is the set of evidence sources in \(R_H\).
The algorithm creates buckets
\[
    B_0,B_1,\ldots,B_D.
\]

It also maintains an arrival list \(A(v)\) for each node that is reached by a
selected path token.  Each arrival is either \(\MinMode\) or \(\MaxMode\).
An arrival represents one still-distinct evidence component, which can be
either a single evidence source or a group of sources that merged earlier.

The local operation
\[
    \operatorname{AddArrival}(v,s)
\]
performs the following steps:

\begin{enumerate}
    \item If this is the first arrival at \(v\), create \(A(v)\) and schedule
          \(v\) exactly once by appending it to bucket \(B_{d(v)}\).
    \item Append mode \(s\) to \(A(v)\).
\end{enumerate}

Later arrivals append to the same list but do not schedule the node again.

\subsection{Step 8: Seed One Arrival per Evidence Node}

For each relevant evidence node \(s\in\mathcal{S}_H\), initialize
\[
    \operatorname{AddArrival}\bigl(s,\Strongest(s)\bigr).
\]

At this point no complete path has been materialized.  Each evidence source has
only a token that records which precomputed extreme it initially prefers.

\subsection{Step 9: Sweep Buckets from Farthest to Nearest}

Buckets are processed in the order
\[
    B_D,B_{D-1},\ldots,B_0.
\]

For every scheduled node \(v\):

\begin{enumerate}
    \item Add \(v\) to the selected-node set.
    \item If \(v=H\), stop processing this token.
    \item Inspect the number of arrivals in \(A(v)\).
    \item If \(|A(v)|=1\), preserve that arrival's current min/max mode.
    \item If \(|A(v)|>1\), treat \(v\) as a merge and reset the continuation
          mode to \(\Strongest(v)\).
    \item Select \(p_{\min}(v)\) or \(p_{\max}(v)\), according to the chosen
          mode.
    \item Add that edge and its destination to the selected sets.
    \item Emit exactly one arrival at the destination, carrying the chosen
          mode.
\end{enumerate}

Emitting one arrival after a merge collapses all incoming evidence components
into one component for the remainder of their shared suffix.

\subsection{Why the Distance Order Is Correct}

For every relevant edge \((v,u)\), the longest-hop recurrence guarantees
\[
    d(v) \geq d(u)+1,
\]
and hence
\[
    d(v)>d(u).
\]

Every emitted arrival therefore moves to a strictly lower, not-yet-processed
bucket.  By induction, all possible arrivals at a node have been collected
before that node's bucket is processed.  This remains true even if an edge
skips one or more distance values.

This ordering provides three useful guarantees:

\begin{enumerate}
    \item a node is scheduled and processed at most once;
    \item a merge sees all of its incoming evidence components;
    \item selecting a previously unused suffix can safely create a new merge
          nearer to \(H\), because the newly reached node has not yet been
          processed.
\end{enumerate}

\subsection{How Nested Merges Are Handled}

Suppose two paths merge at \(M\), and the resulting shared path later meets a
third evidence path at \(N\):

\begin{verbatim}
E1 --> A --\
           M --> C --\
E2 --> B --/          N --> ... --> H
                    /
E3 --> D -----------/
\end{verbatim}

At \(M\), the arrivals from \(E_1\) and \(E_2\) collapse into one token.  All
edges in their distinct prefixes have already been selected and remain in the
output.  The collapsed token moves from \(M\) toward \(N\).
The suffix chosen at \(M\) may be one that neither incoming token would have
followed without the merge.  Its call to \(\operatorname{AddArrival}\) still
records the new token at \(N\) before the lower-distance bucket for \(N\) is
processed.

At \(N\), the algorithm sees two arrivals:

\begin{enumerate}
    \item one token representing the already-merged component
          \(\{E_1,E_2\}\);
    \item one token representing \(E_3\).
\end{enumerate}

It therefore performs another merge, preserves all three prefixes, and emits
one token along a single continuation from \(N\) to \(H\).

\subsection{Worked Pruning Example}

Consider the following path-ranking weights:

\begin{center}
\begin{tabular}{cl@{\qquad}cl}
\toprule
Edge & \(r_e\) & Edge & \(r_e\) \\
\midrule
\(E_1\to A\) & \(2\)       & \(E_2\to B\) & \(0.5\) \\
\(A\to M\)   & \(2\)       & \(B\to M\)   & \(0.5\) \\
\(M\to C\)   & \(4\)       & \(M\to X\)   & \(0.25\) \\
\(C\to N\)   & \(1\)       & \(X\to N\)   & \(1\) \\
\(E_3\to D\) & \(3\)       & \(D\to N\)   & \(1\) \\
\(N\to H\)   & \(2\)       &                 & \\
\bottomrule
\end{tabular}
\end{center}

For \(E_1\), the path through \(C\) has ratio
\[
    2\cdot 2\cdot 4\cdot 1\cdot 2=32,
\]
whereas the path through \(X\) has ratio \(2\).  Thus \(E_1\) initially
selects the maximum mode and prefers \(C\).

For \(E_2\), the path through \(C\) has ratio
\[
    0.5\cdot 0.5\cdot 4\cdot 1\cdot 2=2,
\]
whereas the path through \(X\) has ratio
\[
    0.5\cdot 0.5\cdot 0.25\cdot 1\cdot 2=0.125.
\]
Because
\[
    |\log 0.125|>|\log 2|,
\]
\(E_2\) initially selects the minimum mode and prefers \(X\).

Both tokens reach \(M\).  At \(M\), the complete suffix ratios are
\[
    M\to C\to N\to H: 4\cdot 1\cdot 2=8,
\]
and
\[
    M\to X\to N\to H: 0.25\cdot 1\cdot 2=0.5.
\]
Since \(|\log 8|>|\log 0.5|\), the merge selects the maximum suffix through
\(C\).  The unused \(M\to X\) suffix is never added to the result.

The merged token then reaches \(N\), where it joins the token from \(E_3\).
The output retains
\[
\begin{split}
    &E_1\to A\to M,\\
    &E_2\to B\to M,\\
    &M\to C\to N,\\
    &E_3\to D\to N,\\
    &N\to H,
\end{split}
\]
but excludes \(M\to X\to N\).

\subsection{Step 10: Construct the Returned Graph}

After the sweep, the algorithm scans the original graph and copies only
selected nodes and edges into a new graph.  Graph metadata is preserved, and
nodes, edge objects, tag collections, and evidence details are cloned so that
mutating the pruned result does not mutate the input graph.

If no relevant evidence reaches \(H\), the returned graph contains \(H\) and
no edges.

\subsection{Pruning Invariants}

The bucket sweep establishes the following invariants.

\begin{enumerate}
    \item \textbf{Reachability:} every selected node has a selected path to
          \(H\).
    \item \textbf{Prefix preservation:} every edge selected before a merge
          remains in the returned graph.
    \item \textbf{Single continuation:} every selected node other than \(H\)
          has one selected outgoing edge toward \(H\).
    \item \textbf{Complete merge detection:} every component that reaches a
          node arrives before that node is processed.
    \item \textbf{Nested-merge safety:} a merged component can participate in
          later merges without reopening or copying its earlier prefixes.
\end{enumerate}

\subsection{Pruning Complexity}

Let \(V_H=|R_H|\), let \(E_R\) be the number of target-relevant edges, and let
\(K=|\mathcal{S}_H|\).

\begin{center}
\begin{tabular}{ll}
\toprule
Operation & Expected cost \\
\midrule
Build graph indexes & \(O(|V|+|E|)\) \\
Find \(R_H\) & \(O(V_H+E_R)\) \\
Topological min/max DP & \(O(V_H+E_R)\) \\
Arrival/bucket sweep & \(O(V_H+K)\) \\
Clone selected output & \(O(|V|+|E|)\) \\
\bottomrule
\end{tabular}
\end{center}

With expected constant-time dictionary and set operations, the complete
implementation is therefore
\[
    O(|V|+|E|+K),
\]
and hence linear because \(K\leq |V|\).  The implementation achieves this by
reusing one DP table and never rerunning a graph search for each merge.
Auxiliary storage is linear in the graph and the number of arrivals.

\section{Stage II: Bayes-Factor Calculation}

\subsection{Input and Structural Requirement}

Write the already-pruned graph as
\[
    G_H^{\star}=(V_H^{\star},E_H^{\star}).
\]
The calculator receives

\begin{enumerate}
    \item an already-pruned graph \(G_H^{\star}\);
    \item the hypothesis identifier \(H\);
    \item a dictionary \(\mathcal{B}_{\mathrm{leaf}}\) of positive leaf Bayes
          factors.
\end{enumerate}

The relevant portion of the input must be a DAG in which every relevant
non-hypothesis node has at most one continuation toward \(H\), and \(H\) has no
continuation within the relevant subgraph.  Multiple child branches may enter
the same parent; those are precisely the contributions multiplied at that
parent.

\subsection{Bayes-Factor Recurrence}

For a node \(v\), let
\[
    N(v)
      = \{u : (u,v)\in E_H^{\star}\}
\]
be its direct children in the pruned graph.

The node Bayes factor is
\[
    \BF_v
      =
      \begin{cases}
        \mathcal{B}_{\mathrm{leaf}}[v],
            & N(v)=\varnothing,\\[0.8em]
        \displaystyle
        \prod_{u\in N(v)} \BF_{u\rightarrow v},
            & N(v)\neq\varnothing.
      \end{cases}
\]

For edge \(e=(u,v)\), transform the child's Bayes factor through the edge:
\[
    \boxed{
    \BF_{u\rightarrow v}
      =
      \frac{
        \BF_u a_e + (1-a_e)
      }{
        \BF_u b_e + (1-b_e)
      }}.
\]

The final return value is \(\BF_H\).

\subsection{Interpretation of the Edge Transform}

The transform maps the strength of the entire child subtree into a Bayes-factor
contribution for its parent.  It has several useful boundary cases:

\begin{itemize}
    \item If \(a_e=b_e\), then \(\BF_{u\rightarrow v}=1\) for every positive
          \(\BF_u\).  The edge is BF-neutral.
    \item If \(a_e=1\) and \(b_e=0\), then
          \(\BF_{u\rightarrow v}=\BF_u\).
    \item If \(a_e=0\) and \(b_e=1\), then
          \(\BF_{u\rightarrow v}=1/\BF_u\).
    \item If \(\BF_u=1\), then \(\BF_{u\rightarrow v}=1\), regardless of
          the valid edge probabilities.
\end{itemize}

Values with \(a_e<b_e\) are valid; depending on the child BF, such an edge can
reverse the direction of the child's contribution.  No sign is inferred from
the edge's \texttt{support} or \texttt{rebut} label.

The schema default \(a_e=b_e=0.5\) is therefore neutral.  Existing or seeded
edges contribute no BF effect until their probabilities are calibrated.

\subsection{Step 1: Restrict Calculation to the Hypothesis Subgraph}

As in pruning, the calculator begins at \(H\) and follows child edges backward
to collect all nodes that can reach \(H\).  Disconnected nodes and edges are
not evaluated.

\subsection{Step 2: Count Unresolved Child Contributions}

For every relevant node \(v\), initialize
\[
    q_{\mathrm{child}}(v)=|N(v)|
\]
and a multiplicative accumulator
\[
    C(v)=1.
\]

Every node with \(q_{\mathrm{child}}(v)=0\) is a leaf and is placed in the
ready queue.

\subsection{Step 3: Evaluate Leaves}

When a leaf \(v\) leaves the queue, retrieve
\[
    \BF_v=\mathcal{B}_{\mathrm{leaf}}[v].
\]

The value must exist and must be strictly positive.  Extra dictionary values
for irrelevant or non-leaf nodes are not needed by the recurrence.

\subsection{Step 4: Evaluate Internal Nodes}

An internal node is queued only after every direct child contribution has been
accumulated.  Its Bayes factor is therefore already available as
\[
    \BF_v=C(v).
\]

\subsection{Step 5: Transform and Propagate to the Parent}

For every relevant continuation edge \(e=(v,u)\):

\begin{enumerate}
    \item validate \(a_e,b_e\in[0,1]\);
    \item calculate \(\BF_{v\rightarrow u}\) with the edge transform;
    \item update
          \[
              C(u)\leftarrow C(u)\,\BF_{v\rightarrow u};
          \]
    \item decrement \(q_{\mathrm{child}}(u)\);
    \item if \(q_{\mathrm{child}}(u)=0\), enqueue \(u\).
\end{enumerate}

The arithmetic is checked decimal arithmetic in the implementation, so an
overflow is reported rather than silently wrapping.

\subsection{Step 6: Detect Cycles and Validate Compatibility}

Let \(n_{\mathrm{processed}}\) be the number of nodes removed from the ready
queue.  If
\[
    n_{\mathrm{processed}}<|R_H|,
\]
then some relevant nodes never reached zero unresolved children, which implies
a cycle.

The calculator also rejects an unpruned relevant node with multiple
continuations toward \(H\).  This prevents a shared child subtree from being
propagated along competing parent paths and counted more than once.

\subsection{Step 7: Return the Hypothesis Bayes Factor}

After all relevant nodes have been evaluated, return the memoized value
\[
    \BF_H.
\]

No node odds are read, divided, multiplied, or updated during this process.

\subsection{Worked Bayes-Factor Example}

Consider the pruned graph
\[
    E_1\longrightarrow M\longrightarrow H,
    \qquad
    E_2\longrightarrow M,
\]
with leaf values
\[
    \BF_{E_1}=4,
    \qquad
    \BF_{E_2}=2.
\]

Let edge \(E_1\to M\) have
\[
    a_1=\frac{3}{4},
    \qquad
    b_1=\frac{1}{4}.
\]
Its contribution is
\[
    \BF_{E_1\rightarrow M}
      =
      \frac{4(3/4)+(1-3/4)}{4(1/4)+(1-1/4)}
      =
      \frac{13/4}{7/4}
      =
      \frac{13}{7}.
\]

Let edge \(E_2\to M\) have \(a_2=0\) and \(b_2=1\).  Then
\[
    \BF_{E_2\rightarrow M}
      =
      \frac{2(0)+(1-0)}{2(1)+(1-1)}
      =\frac{1}{2}.
\]

The two independent child contributions multiply at \(M\):
\[
    \BF_M
      =\frac{13}{7}\cdot\frac{1}{2}
      =\frac{13}{14}.
\]

Finally, let edge \(M\to H\) have
\[
    a_3=\frac{4}{5},
    \qquad
    b_3=\frac{1}{5}.
\]
Then
\begin{align*}
    \BF_H
      &=
      \frac{(13/14)(4/5)+(1-4/5)}
           {(13/14)(1/5)+(1-1/5)}\\[0.4em]
      &=
      \frac{66/70}{69/70}
       =\frac{22}{23}
       \approx 0.95652.
\end{align*}

This value is returned directly.  The calculator does not apply it to any
prior odds.

\subsection{Bayes-Factor Complexity}

Every relevant node is queued once, and every relevant edge is transformed
once.  The target-local traversal and recurrence therefore cost
\[
    O(|V_H^{\star}|+|E_H^{\star}|),
\]
with target-local auxiliary structures of
\[
    O(|V_H^{\star}|+|E_H^{\star}|).
\]
The public call first constructs adjacency indexes for the entire input graph,
so its total time and storage bounds are
\[
    O(|V|+|E|).
\]

\section{End-to-End Properties}

The two stages together provide the following behavior.

\begin{enumerate}
    \item Pruning removes competing continuations before BF multiplication.
    \item Evidence-specific prefixes survive merges.
    \item Every merged component has one shared suffix to \(H\).
    \item Each retained edge's two probabilities are preserved when the graph
          is cloned.
    \item The BF calculator visits each retained node and edge once.
    \item Prior and posterior odds remain unchanged.
\end{enumerate}

\section{Validation and Failure Conditions}

The implementation rejects or reports the following conditions:

\begin{center}
\begin{tabular}{p{0.37\textwidth}p{0.53\textwidth}}
\toprule
Condition & Result \\
\midrule
Missing hypothesis & calculation or pruning stops \\
Duplicate node or edge identifier & pruning stops \\
Dangling edge endpoint & graph-context construction stops \\
Relevant cycle & topological processing stops \\
Nonpositive relevant \(r_e\) & pruning stops \\
Missing or nonpositive leaf BF & BF calculation stops \\
\(a_e\) or \(b_e\) outside \([0,1]\) & BF calculation stops \\
Multiple continuations in BF input & caller is instructed to prune first \\
Cancellation requested & the active operation stops \\
Decimal overflow & checked arithmetic throws \\
\bottomrule
\end{tabular}
\end{center}

\section{Implementation Choices Requiring Product Confirmation}

The current implementation makes several explicit choices where the general
problem statement admits more than one interpretation.

\subsection{Path Metric}

Pruning uses \texttt{ImportanceToParent} as a multiplicative path likelihood
ratio.  It does not derive a pruning score from \(a_e\) and \(b_e\).  Changing
this requires defining how the nonlinear BF transform should be scored before
the child Bayes factor is known.

\subsection{Meaning of ``Strongest''}

The strongest path is the path farthest from neutral BF \(1\):
\[
    \max |\log(\text{path LR})|.
\]
This treats strong supporting and strong counter paths symmetrically.

\subsection{Merge Objective}

At a merge, the implementation chooses the merge node's locally strongest
suffix.  This is a deterministic greedy structural rule.  It is not claimed
to optimize a global objective such as the sum or product of all evidence-path
strengths, because no such aggregate objective is presently defined.

\subsection{Tie Rules}

An exact absolute-strength tie selects the maximum path.  Equal computed DP
scores use ordinal edge identifier and then destination identifier as stable
tie-breaks.  The equality is equality of the computed log values, not symbolic
equality of the original products.

\subsection{Evidence Kinds}

Both \texttt{evidence} and \texttt{objection} nodes start pruning tokens.  If
the domain later introduces other evidence-like kinds, that policy should be
centralized or replaced with an explicit evidence-node input set.

\subsection{Posterior-Odds Integration}

Posterior integration remains intentionally out of scope.  Before adding it,
the application must decide whether stored values are ordinary odds or
log-odds:

\begin{itemize}
    \item for ordinary odds,
          \(O_{\mathrm{post}}=O_{\mathrm{prior}}\BF_H\);
    \item for log-odds,
          \(\log O_{\mathrm{post}}
             =\log O_{\mathrm{prior}}+\log\BF_H\).
\end{itemize}

That decision should be implemented in a separate integration layer rather
than inside either algorithm described here.

\section{Mapping from the Document to the C\# Implementation}

\begin{center}
\begin{tabular}{p{0.36\textwidth}p{0.54\textwidth}}
\toprule
Document operation & Implementation member \\
\midrule
Collect \(R_H\) & \texttt{CollectNodesThatReachHypothesis} \\
Calculate \(m,M,p_{\min},p_{\max},d\) & \texttt{CalculatePathMetrics} and
    \texttt{CalculateNodePathMetrics} \\
Choose min/max extreme & \texttt{SelectStrongest} \\
Register and schedule a token & local function \texttt{AddArrival} \\
Perform compatible merge sweep & \texttt{SelectCompatiblePaths} \\
Clone the selected result & \texttt{CreatePrunedGraph} \\
Retrieve a leaf BF & \texttt{GetLeafBayesFactor} \\
Apply edge recurrence & \texttt{TransformThroughEdge} \\
Calculate \(\BF_H\) & \texttt{GraphBayesFactorCalculator.Calculate} \\
\bottomrule
\end{tabular}
\end{center}

\section{Summary}

The pruning algorithm first computes the minimum and maximum log-LR path from
every target-relevant node to \(H\).  It then propagates evidence-component
tokens through longest-distance buckets.  A single token preserves its chosen
extreme, while multiple tokens at one node trigger a merge and one locally
strongest shared suffix.  This preserves all incoming prefixes and naturally
handles later nested merges.

The BF calculator then processes the pruned DAG from leaves toward \(H\).
Explicit leaf Bayes factors are transformed through edge probabilities, direct
child contributions are multiplied, and the resulting \(\BF_H\) is returned.
Posterior odds are neither consulted nor modified.

\end{document}
